% =============================================================
% Mental Model
% Chapter order is set in ../main.tex; the rendered chapter number
% comes from \chapter{} below.
% =============================================================

\chapter{Mental Model}
\label{ch:mental}

\todo{Chapter scaffold. Each section has a one- to three-paragraph
stub indicating intended content. Fleshing out is deferred to a
later writing pass.}

This chapter develops the working mental model that the rest of the
book takes for granted. The angle is deliberately narrow: not how a
transformer is trained, not why attention is a useful operation,
not the historical chain that runs from RNNs to GPT-5. The reader
is assumed to be an engineer who will treat the language model as
a service with a calling convention and who needs to understand
the model's behavior --- its strengths, its limits, the cost
structure of using it --- well enough to design systems that hold
together. We open with the operational picture, work through the
mechanics that follow from it, and close with a taxonomy of
characteristic failures.

\section{What a Language Model Is}
\label{sec:mental-what}

A language model, for the purposes of this book, is a function. The
argument is a sequence of tokens; the value is a probability
distribution over what the next token might be. That is the entire
interface. Everything the engineer does with a model --- chat,
tool calling, retrieval-augmented generation, agentic loops --- is
built on top of this single shape.

\begin{definitionbox}
A \textbf{language model} is a stateless function from a sequence
of tokens to a probability distribution over a fixed vocabulary
of \emph{next} tokens. \emph{Memory} is the text the host hands
the model on the current call; there is no persistent state across
calls. \emph{Generation} is autoregressive: sample a token from
the distribution, append it to the input, repeat.
\end{definitionbox}

\noindent
Two consequences follow from naming the model ``a function'' so
deliberately, and they are worth dwelling on because the rest of
the book leans on both.

The first: a function does not remember. There is no hidden state
carried across calls, no internal variable that knows what the
user said five minutes ago, no ``memory'' in any conventional
sense. The model that processes our second HTTP request is the
same model that processed our first, with no recollection that the
first one happened. Anything the model appears to ``know'' about
the conversation must be in the input it sees on the
\emph{current} call, because there is nowhere else for that
knowledge to live.

The second: a function with a probabilistic output is not, by
itself, a text generator. The model gives us a distribution over
next tokens; turning a distribution into a single next token
requires \emph{sampling}. The runtime samples one token, appends
it to the input, sends the whole thing back to the model, samples
again, repeats. This is autoregressive generation: each output
token feeds back into the input that produces the next one.
Temperature, top-$p$, top-$k$, and the various choices around
them (\S\ref{sec:mental-sampling}) are the knobs on this loop.

\begin{examplebox}
\small
\textbf{One step of the loop.} Given the input
``\code{The capital of Switzerland is}'', a Gemma-class model
returns a distribution dominated by the token \code{ Bern}, with
smaller probability mass on \code{ Berne}, \code{ Zurich},
\code{ the}, and a long tail of less-likely candidates. Greedy
sampling at temperature zero picks \code{ Bern}. The host appends
it to the input, producing
``\code{The capital of Switzerland is Bern}'', and calls the model
again. The next-token distribution this time is dominated by
\code{.}; greedy sampling picks it; the input is now
``\code{The capital of Switzerland is Bern.}''. And so on, until
the model emits an end-of-sequence token or the runtime stops it.
\end{examplebox}

\noindent
Putting these two facts together: \emph{a language model is a
stateless conditional sampler over a token vocabulary.} That
mouthful is worth internalizing, because it explains every
``memory'' trick the rest of the book will introduce, and every
restriction that follows from one.

The chat history we hold across multi-turn conversations is text
the host appends to the input on every turn. The system prompt
that gives the agent its persona is text at the start of the input
on every turn. The tool advertisements that tell the model about
available functions (Chapter~\ref{ch:toolcalls}) are text in the
input on every turn. The retrieved passages of a RAG system
(Chapter~\ref{ch:context}) are text in the input on every turn.
The model's own previous responses, the user's earlier questions,
the documentation pulled in from a vector store, the reasoning
steps we asked the model to ``think through'' out loud --- all of
it is text in the input on every turn. There is no hierarchy of
importance the model implicitly understands; there is one big
string, lightly segmented by role tags (\S\ref{sec:work-chats}) so
the model can tell which part is which.

Once this picture is in place, a great deal of agent design
becomes legible:

\begin{itemize}[leftmargin=1.5em,itemsep=0.2em]
\item The \emph{context window} of a model is the maximum length
  of this input string. Every byte of system prompt, history,
  tool catalog, and retrieved context counts toward the same
  budget (\S\ref{sec:mental-context}).
\item ``Forgetting'' an old turn is, mechanically, just dropping
  it from the input list before the next call. The model has no
  perspective on what was there before; nothing complains.
\item ``Editing'' the model's own previous response is rewriting
  that part of the input. The model on the next call accepts the
  rewrite as ground truth. Agent harnesses use this to compress
  intermediate results.
\item Agent harnesses with explicit ``memory'' subsystems
  (Chapter~\ref{ch:context}) are deciding, on every turn, what
  to put back in this input.
\end{itemize}

\noindent
This single fact --- the input on this call is the only place
memory lives --- is the most useful piece of mental scaffolding
for everything that follows. We will return to it, sometimes
explicitly, on every page where a design decision is made about
what goes into the model's input.

\subsection{Model types}
\label{sec:mental-types}

Language models differ along several axes that matter for system
design without changing the underlying mental model. The reader
will encounter these dimensions repeatedly throughout the book;
three are consequential enough to name here. We treat the
operational question of which to pick for a given task in
Chapter~\ref{ch:working}; this section just maps the territory.

\paragraph{Multimodal models.}
The interface generalizes naturally beyond text. The input is no
longer a sequence of text tokens but a sequence of \emph{content
blocks}, each typed --- text, image, audio, occasionally video.
The model attends to all of them inside the same context window
and produces text out the other end. By 2026, vision input is
standard across the frontier (OpenAI's GPT family, Anthropic's
Claude, Google's Gemini, Meta's Llama, Google's Gemma); audio
input is standard on most of the same; native image output
remains less common and is more often handled by a separate
specialist model invoked as a tool call
(\S\ref{sec:taxonomy}). The agent-design implication is that
several tools which used to be necessary --- OCR for screenshots,
transcription for voice notes, layout extraction for PDFs --- are
subsumed by a multimodal model directly. The mental model is
unchanged: a stateless conditional sampler over a token
vocabulary, where some of those tokens happen to encode pixels or
audio frames rather than letters.

\paragraph{Open-weights and local models.}
``Open weights'' means the trained parameters of the model are
publicly downloadable, typically under a license that permits
redistribution and fine-tuning. Crucially, ``open weights'' does
\emph{not} mean ``open source'': the training data and training
code are usually not released, and the reader who wants to
reproduce the training run typically cannot. What they
\emph{can} do is download the weights, run them on their own
hardware, fine-tune them on their own data, host them on private
infrastructure, and quantize them for cheaper inference. Each of
those is impossible with a closed-weight hosted model, and each
matters for some class of practical use --- private data,
reproducibility (\S\ref{sec:intro-approach}), regulatory
constraints, latency budgets, cost ceilings. The trade-off is
capability: as of 2026 the strongest open-weights models trail
the strongest closed-weight ones on the hardest benchmarks, but
the gap closes every quarter and for many practical workloads
the open-weights option is competitive. Major open-weights
families to be aware of: Meta's Llama, Alibaba's Qwen, DeepSeek,
Mistral, Microsoft's Phi, Google's Gemma --- the model this book
uses by default. Chapter~\ref{ch:working} treats the choosing-a-model
question in detail; the relevant point here is that ``open
weights'' is a real and increasingly viable design choice, not a
hobbyist concession.

\paragraph{Reasoning / thinking models.}
A relatively recent variant in which the model is trained to
spend extra tokens on a private chain of thought \emph{before}
committing to its final answer. The mechanism is the same
autoregressive loop described above, but the convention is that
some prefix of the output --- often delimited by special tokens,
sometimes hidden from the API consumer --- is ``thinking'' rather
than the user-facing answer; only the post-thinking text is
returned. The trade-off is straightforward. On tasks that benefit
from multi-step reasoning --- mathematics, planning, code that
requires careful tracing, formal verification --- reasoning
models are noticeably stronger, often by a wide margin. On tasks
that don't --- lookup, summarization, simple extraction, format
conversion --- they are merely slower and more expensive, since
the thinking tokens are still tokens the user pays for. Major
families that ship a thinking mode as of 2026 include OpenAI's
o-series, Anthropic's ``extended thinking'', Google's Gemini
``thinking'' variants, DeepSeek-R1, and Qwen-QwQ. We treat
reasoning as one axis in choosing a model
(\S\ref{sec:work-reasoning}), not a separate species; the mental
model of \S\ref{sec:mental-what} carries over unchanged.

\section{Tokens and Tokenization}
\label{sec:mental-tokens}

The model does not see characters and does not see words; it sees
\emph{tokens}. A token is the unit of input the model is trained
on: an integer index into a fixed-size vocabulary, mapped at the
very first layer of the network to a learned vector in the
model's embedding space. Everything downstream --- the attention
pattern, the layer arithmetic, the next-token prediction ---
operates on these vectors. The character ``r'' has no first-class
representation inside the model; its existence is implicit in
whichever tokens happen to contain it, and only to the extent
that the training corpus correlated those tokens with
character-level identity in some downstream task.

\subsection{What to train on}
\label{sec:mental-tokens-what}

The vocabulary cannot grow without bound. Each entry costs an
embedding vector and a row in the output projection, which means
both memory and parameters; the model must be told, before
training begins, that it has exactly $N$ possible tokens to
predict. Two naive choices for what those $N$ tokens should be
present themselves immediately, and both fail in instructive
ways.

\paragraph{Words.}
Tempting, because language is composed of words and a word-level
model would be the most efficient at representing common text:
a sentence of ten words would be a sequence of ten tokens. The
trouble is that the world does not helpfully restrict itself to
a fixed vocabulary. Names, brand identifiers, scientific terms,
neologisms, internet slang, code identifiers, large numbers,
foreign-language fragments, typos --- every one of these
produces ``words'' the model has never seen. A purely word-based
vocabulary either grows unboundedly to chase them all, or falls
back to a single \code{<UNK>} token that collapses every unknown
into the same vector. Neither option is acceptable. Real text
leaks through the cracks of any finite word list, and the model
must have something to say about \emph{strawpberry},
\code{f1bcd2}, \emph{Belfanti}, and \code{quaternion} even when
none of those have appeared in training.

\paragraph{Letters or bytes.}
The opposite extreme. There are only 256 possible byte values,
which fits in a single byte of indexing and admits any input
string without exception. Every hash, every typo, every unknown
name decomposes into bytes the model has seen many times. The
trouble here is the other side of the same coin: the model must
now learn from scratch that \code{t}, \code{h}, \code{e}
together mean \emph{the}; that \emph{the} appears at the start
of countless English sentences; that the bigram \code{th} is
much more common than \code{tg}. The most useful patterns in
language are not at the byte level, they are several bytes long,
and a byte-level model burns much of its capacity reconstructing
those patterns at every position. Sequences also become very
long: a sentence of ten English words might be sixty bytes,
six times the work for the same content.

\subsection{The compromise: bin what's frequent}
\label{sec:mental-tokens-compromise}

The middle ground is to let the data decide. Train the
vocabulary on the same corpus the model will eventually be
trained on; give frequent character sequences their own token,
and reserve single-byte tokens as a fallback for the rare cases.
The resulting vocabulary mixes granularities: \code{the} is one
token because it occurs everywhere; \code{ing} and \code{tion}
are tokens because they recur as suffixes; rare words break into
several pieces; entirely unfamiliar input falls all the way back
to bytes.

The dominant algorithm for producing such a vocabulary is
Byte-Pair Encoding (BPE)~\citep{sennrich2016bpe}, with close
relatives WordPiece and Unigram in production use. Conceptually
the recipe is straightforward.

\begin{enumerate}[label=(\arabic*),leftmargin=2em,itemsep=0.2em]
\item Start with a vocabulary of all single bytes. Every input
  string is representable as a sequence of these.
\item Count, across the training corpus, the frequency of every
  pair of adjacent tokens.
\item Merge the most-common pair into a single new token, added
  to the vocabulary. Re-tokenize the corpus.
\item Repeat until the vocabulary reaches the target size.
\end{enumerate}

\noindent
A run that has seen a great deal of English will discover, very
early, that \code{t} followed by \code{h} is the most common
pair, and merge it into \code{th}; soon after, that \code{th}
followed by \code{e} is the most common pair, and merge into
\code{the}; soon after, that \code{ ing} (with its leading
space) is worth its own token. Run for long enough on a corpus
large enough, and the vocabulary contains common English words
as single tokens, inflectional fragments (\code{ing},
\code{tion}, \code{ly}), most individual letters, and almost
everything in between. Modern vocabularies sit somewhere between
$\sim$50\,000 entries (early GPT-3, Llama~2) and $\sim$256\,000
(Gemma~3 and~4, GPT-4o). At inference time, the tokenizer that
converts an input string into a token sequence is a
longest-match algorithm against this vocabulary.

\subsection{Practical consequences}
\label{sec:mental-tokens-consequences}

Three implications follow from this design and recur throughout
the book.

\begin{itemize}[leftmargin=1.5em,itemsep=0.4em]

\item \textbf{Token counts are not character counts and are not
  word counts.} A useful rule of thumb is that 1000~tokens is
  roughly 750~English words, or about 4000 characters of running
  English prose. The ratio swings widely: code is denser than
  prose because indentation, punctuation, and operators each
  count; structured data (JSON, XML) denser still; rare technical
  vocabulary denser again. A screen of Python is two to three
  thousand tokens; the same screen of plain English is twelve to
  fifteen hundred.

\item \textbf{The cost-per-character is not language-neutral.}
  BPE training corpora have been overwhelmingly English. Latin
  scripts share a tokenizer-friendly structure with English and
  pay a small premium. Cyrillic, Arabic, Devanagari, and CJK
  scripts pay much more: a Mandarin character can be two or three
  tokens, because rare characters fall back to byte-level encoding
  within the tokenizer. On per-token-priced APIs this is a
  multiplicative cost on every interaction in a non-English
  language, and on every prompt that carries non-English text.

\item \textbf{Different models use different tokenizers.} Gemma~4
  ships its own SentencePiece-trained vocabulary of
  $\sim$256\,000 entries; OpenAI's GPT-4 family uses
  \code{cl100k\_base}; GPT-4o uses \code{o200k\_base}; Anthropic
  publishes a \code{count\_tokens} endpoint rather than the
  tokenizer itself. The same string can take 500~tokens on one
  model and 600~on another. Counting accurately --- whether to
  budget context, estimate cost, or fit a passage under a length
  limit --- requires the tokenizer the target model is actually
  using.

\end{itemize}

\noindent
The mechanism is most readable in the small. The example below
uses OpenAI's \code{o200k\_base} (the encoding behind GPT-4o)
because it is freely available; the qualitative pattern is the
same across modern tokenizers.

\begin{examplebox}
\small
\textbf{What the tokenizer actually emits.} Each input below is
followed by its token count and the literal sequence of token
strings that the tokenizer produced.

\begin{itemize}[leftmargin=1.5em,itemsep=0.1em]
\item \code{"hello world"} --- 2 tokens --- \code{["hello", " world"]}
\item \code{"strawberry"} --- 3 tokens --- \code{["st", "raw", "berry"]}
\item \code{"strawpberry"} --- 4 tokens --- \code{["st", "raw", "p", "berry"]}
\item \code{"def count\_letter(word, letter):"} --- 7 tokens --- \code{["def", " count", "\_letter", "(word", ",", " letter", "):"]}
\item \code{"François"} --- 2 tokens --- \code{["Fran", "çois"]}
\item \code{"supercalifragilisticexpialidocious"} --- 11 tokens (rare and broken into syllable-sized pieces)
\item the five-character Mandarin sentence \emph{n\v{i} h\v{a}o,
  sh\`{i}ji\`{e}} (``hello, world'') --- 6 tokens, with the
  comma and the most common characters tokenizing one-to-one and
  the rarer characters falling back to multi-byte sequences
\end{itemize}

\noindent
Note in particular how the misspelling \emph{strawpberry} differs
from the correct \emph{strawberry}: the inserted \texttt{p}
appears as a token of its own because the bigram ``rawp'' is not
in the vocabulary. The model sees the difference between the two
spellings as a four-vector vs three-vector input, not as a
single-character change.
\end{examplebox}

\noindent
The closing observation is the one that motivates the entire
strawberry section of Chapter~\ref{ch:toolcalls}: the model has
no character-level representation of its own input. Any task that
requires reasoning about individual letters --- counting them,
permuting them, checking spelling --- is asking the model to
recover character-level information from token-level vectors,
which it can do only weakly and unreliably. The full unpacking
of this failure mode lives in \S\ref{sec:strawberry}; the short
version here is enough to set up the broader taxonomy of
\S\ref{sec:mental-failures}, where tokenizer-level failures are
one cell in the larger picture.

\section{The Context Window}
\label{sec:mental-context}

\todo{Short. Nominal context vs effective context: the
``lost-in-the-middle'' phenomenon, the empirical observation that
recall on facts placed in the middle of a long context degrades
sharply, and that effective utility caps out well below the
advertised maximum. Forward-reference the
Tool-Calls-and-Context section of Chapter~\ref{ch:toolcalls} for
the agent-side implications. Detail of \emph{what fills} the
context (system prompt, tool catalog, history, tool results)
lives there, not here.}

\section{Sampling and Determinism}
\label{sec:mental-sampling}

\todo{What temperature, top-k, top-p are at the level of intuition.
Greedy decoding at temperature zero. The catch: ``temperature 0''
is not actually deterministic across providers, because of
batching, hardware variation, and silent model updates. This has
direct consequences for testing agents
(Chapter~\ref{ch:reliable}). Mention seed parameters where they
exist and where they do not.}

\section{Embeddings}
\label{sec:mental-embeddings}

\todo{Short. An embedding model maps a piece of text to a
fixed-length vector such that semantically similar texts get
nearby vectors. Practical facts the reader needs: typical
dimensionalities (512, 768, 1024, 1536); cosine similarity / dot
product as the comparison function; dramatically cheaper per call
than chat models; the existence of dedicated embedding models
distinct from chat models. Defer the \emph{use} of embeddings
(vector stores, retrieval, memory, deduplication) to the
agent-context chapter where RAG is treated. Frame here as
``another thing the same family of models can do, with a different
output shape.''}

\section{Strengths and Failure Modes}
\label{sec:mental-failures}

\todo{The substantive section. A small taxonomy of what language
models are good at and what they consistently fail at, with the
mechanism for each failure named. This is the section that pays
for the rest of the chapter --- the engineer who has internalized
it will design systems differently.}

\subsection{What models are good at}
\label{sec:mental-strengths}

\todo{One paragraph. Strong: language production and comprehension;
code; pattern completion; well-known facts that were dense in
training data; transformations and summarizations of given text;
in-context learning from a few examples.}

\subsection{Tokenizer-level failures}
\label{sec:mental-tok-fail}

\todo{Counting characters, manipulating individual letters, exact
arithmetic on long numbers. Forward-reference
\S\ref{sec:strawberry}, which works the canonical example through.}

\subsection{Hallucination}
\label{sec:mental-hallucination}

\todo{The model produces plausible content that is not true. Why
this happens (the training objective rewards plausibility, not
truth). Mitigations: tools that ground the model in real data
(\S\ref{sec:taxonomy}), retrieval, verifiers (the Lean example
from \S\ref{sec:intro-now}).}

\subsection{Sycophancy and instruction drift}
\label{sec:mental-sycophancy}

\todo{The model agrees too readily with the user's framing, even
when the framing is wrong. Related: over long contexts the model
drifts away from the original instruction in the system prompt.
Implications for agent design: do not assume that pushback
survives disagreement; periodically re-anchor the system prompt.}

\subsection{Sensitivity to phrasing and formatting}
\label{sec:mental-sensitivity}

\todo{Small changes to the prompt --- punctuation, ordering of
examples, capitalization --- can produce non-trivially different
outputs. This makes evaluation subtle (Chapter~\ref{ch:reliable})
and prompts brittle. Forward-reference the prompting-craft
chapter (Chapter~\ref{ch:tools}).}

\subsection{Distributional and adversarial fragility}
\label{sec:mental-fragility}

\todo{Models perform well on inputs that resemble their training
distribution and worse on inputs that do not. Adversarial
inputs --- including prompt injection, treated in
Chapter~\ref{ch:reliable} --- exploit this gap.}

% Structured Outputs has been moved to Chapter~\ref{ch:working},
% \S\ref{sec:structured}. It is more naturally a property of how
% one calls a model than of the model itself.
